
\chapter{Quy Định Không Có Não}
\markboth{Quy Định Không Có Não}{Quy Định Không Có Não}

\section{``Tóc Không Được Chạm Cổ Áo --- Kể Cả Nam''}
\begin{truyen}{``Tóc Không Được Chạm Cổ Áo --- Kể Cả Nam''}{``Hair Must Not Touch the Collar, Boys Included''}
\chuhoa{N}{am sinh bị gọi lên} phòng giám thị vì tóc vừa chạm cổ áo. Kiểm tra lại thì tóc dài đúng hai milimét so với giới hạn.
\textit{(A male student is called to the discipline office. His hair just barely touches his collar. Measured exactly: two millimetres over the limit.)}

\studentpair{Thưa thầy, hai milimét thôi. Em cắt rồi sẽ dài lại sau một tuần.}{``Teacher, it is only two millimetres. If I cut it now it will grow back in a week.''}

\teacherpair{Quy định là quy định. Không bàn cãi. Gọi phụ huynh vào.}{``Rules are rules. No discussion. Call your parents in.''}

Phụ huynh bỏ buổi sáng lái xe vào trường. Nhìn tóc con. Nhìn giám thị. Không nói gì.
\textit{(The parent takes off a morning to drive to school. Looks at the son's hair. Looks at the discipline teacher. Says nothing.)}

\begin{baihoc}
Quy định tóc tai có thể có lý do chính đáng về tác phong. Nhưng khi thực thi quy định mất nhiều tài nguyên hơn giá trị nó bảo vệ, đó là dấu hiệu quy định đang được dùng như quyền lực, không phải như công cụ giáo dục.
\textit{(Hair regulations may have legitimate reasons around conduct. But when enforcing a rule costs more resources than the value it protects, that is a sign the rule is being used as power, not as an educational tool.)}
\end{baihoc}
\end{truyen}

\section{``Không Được Mặc Áo Lạnh Vào Ngày Không Lạnh Theo Lịch''}
\begin{truyen}{``Không Được Mặc Áo Lạnh Vào Ngày Không Lạnh Theo Lịch''}{``No Jackets on Days That Are Not Officially Cold''}
\chuhoa{T}{rời 18 độ nhưng} theo lịch năm học, đó vẫn là mùa hè. Giám thị nhắc học sinh cởi áo khoác.
\textit{(The temperature is 18 degrees, but according to the school schedule, it is still summer. The discipline teacher tells students to remove their jackets.)}

\studentpair{Thưa thầy, trời lạnh mà.}{``But it is cold, teacher.''}

\teacherpair{Chưa đến ngày được mặc áo lạnh. Cởi ra.}{``It is not yet the approved jacket date. Take it off.''}

\begin{baihoc}
Khi quy định không theo kịp thực tế, người thực thi quy định có hai lựa chọn: dùng lý trí hoặc dùng quyền lực. Bắt học sinh run vì lịch chưa cho phép mặc ấm không phải kỷ luật. Đó là sự cứng nhắc được mặc thêm áo quyền lực.
\textit{(When a rule fails to keep up with reality, the enforcer has two choices: use reason or use authority. Forcing students to shiver because the calendar has not approved warm clothes is not discipline. It is rigidity wearing the costume of authority.)}
\end{baihoc}
\end{truyen}

\section{Họp Lớp Về Quy Định Mà Không Học Sinh Nào Được Góp Ý}
\begin{truyen}{Họp Lớp Về Quy Định Mà Không Học Sinh Nào Được Góp Ý}{A Class Meeting About Rules Where No Student Input Was Taken}
\chuhoa{G}{iáo viên chủ nhiệm:} ``Hôm nay mình họp về nội quy lớp học kỳ mới. Tôi đọc danh sách, các em ký tên xác nhận.''
\textit{(``Today we meet to discuss class rules for the new semester. I will read the list, and you all sign to confirm.'')}

\studentpair{Thưa thầy, tụi em có được góp ý không?}{``Are we allowed to give input, teacher?''}

\teacherpair{Không cần. Quy định đã được nhà trường duyệt.}{``No need. The school has already approved the regulations.''}

\studentpair{Vậy cuộc họp này để làm gì ạ?}{``Then what is the point of this meeting?''}

Câu hỏi đó không được trả lời. Thầy chuyển sang trang kế tiếp.
\textit{(That question goes unanswered. The teacher moves to the next page.)}

\begin{baihoc}
Một cuộc họp lấy ý kiến mà ý kiến không được lấy không phải họp, đó là buổi thông báo có nghi thức giả. Học sinh nhận ra sự khác biệt đó nhanh hơn người lớn tưởng.
\textit{(A consultation meeting that takes no input is not a meeting. It is an announcement with a ceremonial disguise. Students notice this difference faster than adults expect.)}
\end{baihoc}
\end{truyen}

\section{Quy Định Áp Dụng Không Đều}
\begin{truyen}{Quy Định Áp Dụng Không Đều}{Rules Applied Unevenly}
\chuhoa{H}{ai học sinh đến trễ cùng một buổi:} một đứa là lớp trưởng, một đứa là học sinh trung bình.

\textit{(Two students arrive late on the same morning: one is class president, the other is an average student.)}

Lớp trưởng được cho vào lớp không nhắc. Học sinh kia bị yêu cầu đứng ngoài.
\textit{(The class president is let in without comment. The other student is required to stand outside.)}

Cả lớp nhìn thấy. Không ai nói gì. Nhưng tất cả đều hiểu.
\textit{(The entire class sees it. No one says anything. But everyone understands.)}

\begin{baihoc}
Công bằng không phải đối xử tất cả như nhau trong từng tình huống. Nhưng khi sự ưu đãi đến từ vị trí xã hội chứ không phải từ lý do chính đáng, học sinh học được điều không ai muốn dạy: quy tắc là cho kẻ yếu, không phải cho kẻ mạnh.
\textit{(Fairness does not mean treating everyone identically in every situation. But when favoritism comes from social position rather than legitimate reason, students learn something no one intended to teach: rules are for the weak, not the powerful.)}
\end{baihoc}
\end{truyen}

\section{Hình Phạt Tập Thể Vì Một Người}
\begin{truyen}{Hình Phạt Tập Thể Vì Một Người}{Collective Punishment for One Person's Action}
\chuhoa{M}{ột học sinh vẽ bậy lên bàn.} Giáo viên không tìm ra ai. Cả lớp bị phạt ở lại vệ sinh lớp học sau giờ tan.
\textit{(A student draws on the desk. The teacher cannot find out who. The entire class is punished by staying to clean up after school.)}

\studentpair{Nhưng thưa thầy, tụi em không làm.}{``But teacher, we did not do it.''}

\teacherpair{Không ai khai thì cả lớp chịu. Dạy cho tụi em biết bảo nhau.}{``If no one confesses, everyone shares the punishment. It teaches you to look out for each other.''}

Người làm: im lặng vì không ai khai sẽ được lợi. Những người không làm: bực bội, nhưng không dám nói. Logic của hình phạt tập thể tạo ra văn hoá im lặng, không phải văn hoá bảo nhau.
\textit{(The guilty student stays silent because not being caught is in their interest. The innocent students are angry but dare not speak. The logic of collective punishment creates a culture of silence, not a culture of accountability.)}

\begin{baihoc}
Hình phạt tập thể có một sức hấp dẫn về mặt hiệu quả ngắn hạn. Nhưng về dài hạn, nó dạy học sinh rằng công lý là thứ có thể bị phân bổ theo số đông, bất kể ai đúng ai sai.
\textit{(Collective punishment has short-term administrative appeal. But long-term, it teaches students that justice can be distributed by the crowd, regardless of who is right or wrong.)}
\end{baihoc}
\end{truyen}

\begin{baihoc}
Quy định cần có não để tồn tại đúng nghĩa. Không phải quy định nào cũng sai, nhưng quy định nào cũng cần được hỏi: nó đang bảo vệ điều gì, và cái giá phải trả để thực thi nó có xứng với điều đó không.
\textit{(Rules need reason to exist correctly. Not every rule is wrong, but every rule needs to be asked: what is it protecting, and is the cost of enforcing it worth what it protects?)}
\end{baihoc}

\begin{ghinhoanh}
\textbf{Short English to remember:}

A rule without reason is just a wall.

Blind obedience and blind rule-making feed each other.
\end{ghinhoanh}

\begin{tuhocnhanh}
1. Quy định nào trong trường em thấy vô lý nhất? Nó đang bảo vệ điều gì? \textit{(Which school rule do you find most unreasonable? What is it actually protecting?)}

2. Làm thế nào để thay đổi một quy định không hợp lý mà không phá vỡ cả hệ thống? \textit{(How do you change an unreasonable rule without breaking the whole system?)}
\end{tuhocnhanh}

\ngancach
